\begin{table}[htbp]
\centering
\caption{OLS vs 2SLS: Effect of Engineer DC on Infrastructure Outcomes}
\label{tab:iv_main_v2}
\begin{tabular}{lcccc}
\hline\hline
 & \multicolumn{2}{c}{Road completion rate} & \multicolumn{2}{c}{Nightlights per capita} \\
\cmidrule(lr){2-3}\cmidrule(lr){4-5}
 & OLS & 2SLS & OLS & 2SLS \\
 & (1) & (2) & (3) & (4) \\
\hline
Engineer & 0.0281 & 0.0265 & -0.0008 & 0.0386 \\
          & (0.0249) & (0.1993) & (0.0007) & (0.0390) \\[4pt]
First-stage $F$ & --- & 6.94~$\times$ & --- & 6.18~$\times$ \\
Observations & 453 & 453 & 2,339 & 2,339 \\
\hline
District FE  & \checkmark & \checkmark & \checkmark & \checkmark \\
Year FE      & \checkmark & \checkmark & \checkmark & \checkmark \\
\hline\hline
\multicolumn{5}{p{0.90\textwidth}}{\footnotesize
\textit{Notes:} Baseline category: SocialScientist.
Instrument (2SLS): Queue\_refined $=$ opposition\_turnover $\times$
cadre\_queue\_pct, where opposition\_turnover $=1$ when incoming
state government is from a different party than outgoing,
and cadre\_queue\_pct $\in[0,1]$ is officer's percentile position
in the cadre's seniority queue in that year (0 = most senior).
Cols (1)--(2): roads subsample, year $\geq 2001$.
Cols (3)--(4): full panel, 1994--2008.
Clustered SEs at district level. $\checkmark$: $F>10$; $\times$: $F\leq10$.
Stars: *** $p<0.01$, ** $p<0.05$, * $p<0.10$.}
\end{tabular}
\end{table}
